% Chủ sở hữu: TV3 (P5.3a, P5.3b) — Review: TV1
\section{Giới thiệu (Introduction)}
\label{sec:introduction}

Các trò chơi hai người, tổng bằng không, thông tin hoàn hảo và tất định như Tic-Tac-Toe, Caro (Gomoku) hay Cờ tướng từ lâu đã là môi trường chuẩn mực để đánh giá các thuật toán lập kế hoạch và tìm kiếm đối kháng trong Trí tuệ Nhân tạo~\cite{knuth1975analysis,allis1994searching}. Trong khi thuật toán Minimax kết hợp cắt tỉa Alpha-Beta~\cite{knuth1975analysis} khai thác cấu trúc cây trò chơi theo chiều sâu cùng hàm lượng giá heuristic đặc thù miền, thì Tìm kiếm Cây Monte Carlo kết hợp giới hạn tin cậy trên cho cây (MCTS-UCT)~\cite{kocsis2006bandit,browne2012survey} tiếp cận bài toán qua việc lấy mẫu ngẫu nhiên có định hướng dựa trên bài toán máy đánh bạc nhiều tay (Multi-Armed Bandit)~\cite{auer2002finite}. Sự kết hợp giữa tìm kiếm cây và học biểu diễn chính sách/giá trị sau đó đã dẫn tới những bước đột phá lớn như AlphaGo~\cite{silver2016mastering}.

Tuy nhiên, khi chuyển từ không gian trạng thái nhỏ có thể giải vét cạn (Tic-Tac-Toe $3\times 3$) sang không gian trạng thái có hệ số phân nhánh lớn và nhiều bẫy chiến thuật dài hạn (Caro $9\times 9$, $15\times 15$), hiệu năng thực tế của mỗi họ thuật toán phụ thuộc chặt chẽ vào các quyết định thiết kế kỹ thuật: thứ tự duyệt nước đi, hàm lượng giá mẫu, hệ số khám phá $c$, tiêu chuẩn chọn nước đi tại gốc, và chất lượng của chính sách mô phỏng (rollout policy).

\subsection{Câu hỏi Nghiên cứu (Research Questions)}
\label{subsec:rqs}

Để đánh giá định lượng và có kiểm chứng thống kê các yếu tố trên, báo cáo này tập trung trả lời sáu câu hỏi nghiên cứu (RQ1--RQ6) bám sát đặc tả thực nghiệm:

\begin{itemize}
    \item \textbf{RQ1 (Cắt tỉa Alpha-Beta \& Sắp xếp nước trên $3\times 3$ --- R1.1, R1.2, Bảng T1):} Cắt tỉa Alpha-Beta với ba chế độ sắp xếp thứ tự nước đi (ngẫu nhiên, cố định theo ưu tiên tâm--góc--cạnh, và đảo ngược) giúp giảm tỷ lệ nút phải duyệt như thế nào so với Minimax thuần túy mà vẫn bảo toàn tuyệt đối giá trị tối ưu trên mọi thế cờ?
    \item \textbf{RQ2 (Tối ưu hóa Alpha-Beta trên Caro --- R1.3, R1.4, Bảng T2):} Trên bàn Caro $n\times n$, việc kết hợp thu hẹp tập nước ứng viên (\texttt{candidate\_moves}), phát hiện thắng/chặn tức thì, sắp xếp nước theo hàm lượng giá mẫu, tìm kiếm sâu dần (Iterative Deepening) và bảng chuyển vị Zobrist (Transposition Table) cải thiện độ sâu tìm kiếm hiệu dụng và tỷ lệ thắng ra sao trong giới hạn thời gian $\le 1\text{ s}$/nước?
    \item \textbf{RQ3 (Kiểm chứng tính ``Không thua'' của MCTS-UCT trên $3\times 3$ --- R2.8, Bảng T3):} Khi đối đầu với Minimax hoàn hảo trên toàn bộ cây phản ứng qua 20 seed độc lập, MCTS-UCT cần ngân sách mô phỏng tối thiểu $I^*$ bằng bao nhiêu để đạt tỷ lệ thua bằng $0\%$, và những thế cờ bẫy (trap states) nào khiến MCTS thất bại ở ngân sách thấp?
    \item \textbf{RQ4 (Tương tác giữa Hệ số Khám phá $c$ và Ngân sách Mô phỏng --- R3, Hình F1, F2):} Hằng số khám phá $c$ trong công thức UCT tương tác với số lượng mô phỏng $I$ như thế nào đối với tỷ lệ thắng của MCTS-UCT trên cả Tic-Tac-Toe $3\times 3$ và Caro $9\times 9$?
    \item \textbf{RQ5 (Tiêu chuẩn Chọn nước cuối và Biến thể UCT --- R4, B2, B3, Bảng T4, T5):} Ba tiêu chuẩn chọn nước đi cuối tại gốc (Robust Child $\arg\max N$, Max Child $\arg\max W/N$, và Max-Robust) cùng các biến thể công thức UCT khác biệt như thế nào về mặt thống kê và hành vi tại các thế cờ nhạy cảm?
    \item \textbf{RQ6 (Chính sách Rollout L0--L4 và Đối đầu Alpha-Beta vs.\ MCTS --- R5, Bảng T6, T7):} Các chiến lược mô phỏng từ ngẫu nhiên đồng nhất (L0), thắng/chặn tức thì (L1), ưu tiên lân cận (L2), chọn mềm theo mẫu (L3) đến cắt sớm bằng hàm lượng giá (L4) đánh đổi giữa tốc độ mô phỏng (simulations/s) và chất lượng đánh giá như thế nào khi so sánh cùng số iteration và cùng ngân sách thời gian (bao gồm thế cờ bẫy E1), cũng như khi đối đầu trực tiếp với Alpha-Beta trên bàn Caro $15\times 15$?
\end{itemize}
